\definecolor{gray}{rgb}{0.9,0.9,0.9}
\definecolor{lightred}{rgb}{1,0.5,0.5}
\def\a#1{\emph{#1}}
%\def\b#1{\emph{\sout{#1}}}
\def\b#1{\emph{*#1}}
\def\w#1{\emph{\underline{#1}}}

%\def\s#1#2#3#4{
%\parbox[t]{7.8cm}{#2} &
%\parbox[t]{7.8cm}{#4} \smallskip \\ \hline
%}

\def\s#1#2#3#4{
\colorbox{#1}{\parbox[t]{7.8cm}{#2}} &
\colorbox{#3}{\parbox[t]{7.8cm}{#4}} \\ \hline
}

\def\empty#1{
\s{gray}{#1 \textcolor{gray}{X}}{gray}{\textcolor{gray}{X}}
}
\def\emptyx#1#2{
\s{#1}{#2 \textcolor{#1}{X}}{#1}{\textcolor{#1}{X}}
}

\def\title#1{\addcontentsline{toc}{section}{#1}
\emptyx{white}{\emph{#1}}
\s{white}{#1 \textcolor{white}{X}}{white}{\textcolor{white}{X}}
}

%\chapter{Bevezetés}
%
%\section{Hunspell}
%
%A \a{Hunspell} fejlesztése egyéni kezdeményezésből indult
%1997-ben \cite{linuxkonf2002}. A megoldott
%technikai problémákat, és a felhasználásra vonatkozó
%információkat a Magyar Ispell kézikönyv
%\cite{magyarispell} tartalmazza.

\chapter{Szabályok}

A következő táblázat a Magyar Ispell és a Microsoft Office
XP-ben megtalálható helyesírás-ellenőrző magyar helyesírásnak
való megfelelését részletezi.

A vizsgálat során nem minden esetben a Microsoft Office XP,
hanem a Morphologic
cég Helyes-e? programjának Linux operációs rendszerre
2002-ben megjelent változata, az Mspell 0.1 program
szolgáltatta a szükséges információkat, de a tapasztalatok
alapján ez megegyezik a Microsoft Office XP-ben
található Helyes-e? változattal. A Magyar Ispell
képességei a Magyar Ispell szótármodul 0.99.3-as és
a Mispell/Magyar MySpell 0.9.4-es változata alapján
lettek összefoglalva. Az ellenőrző és a szótármodul
ezen változatát a Magyar OpenOffice.org
(\url{http://office.fsf.hu}) közeljövőben megjelenő
változatai tartalmazni fogják.

A táblázat a példaként megadott szóalakokat \a{kurzív}
betűváltozattal jelöli. Az ilyen szóalakok, amennyiben
a magyar helyesírás szerint hibásak, csillaggal bevezetve
jelennek meg (például \b{hejesirás}).
Amennyiben helyesek, de a helyesírás-ellenőrző mégis
hibásnak jelzi ezeket, a
Microsoft Office és az OpenOffice.org szövegszerkesztőkben
ismerős módon, vagyis hullámos aláhúzással
vannak jelölve (például \w{helyesírás}).

A leírások A magyar helyesírás szabályai 11.~kiadása pontjainak megfelelően
számozottak.
Egy-két kiegészítő megjegyzés külön sorban foglal helyett.

%Az egyes leírások háttere színes.
%A szín a kezelésben mutatkozó hibák
%súlyosságát jelzi, bár ez a hiba lehet, hogy
%a helyesen kezelt alakok számához képest kicsi (bár nem
%elhanyagolható) arányban fordul elő.\footnote{Például a Helyes-e?
%egyetlen komoly \a{ly} betűt érintő hibája a
%\b{csevely} elfogadása a \w{csevej}
%tiltása mellett. Ha ehhez még hozzávesszük azt a
%mintegy pár tucat tövet, illetve alapalakot, ahol a jó mellett a hibás alakot
%is elfogadja a Helyes-e? (pl. \a{személyi}--\b{szeméji},
%\a{karvaj}--\b{karvaly} stb.), akkor is csak az
%1\%-ot közelíti meg a hiba az összes szótövet érintő
%\emph{ly}$\rightarrow$\emph{j} tévesztést tekintve. (A Magyar Ispell
%esetében a hiba gyakorlatilag 0\%, eltekintve persze attól,
%hogy egyes szavak, amelyeknek
%mind a \a{j} és az \a{ly} betűs alakváltozatai helyesek,
%a szövegkörnyezettől függően hibásnak bizonyulhatnak: pl. ilyen
%a \a{lakályos}--\a{lakájos}, vagy a \a{gólya}--\a{gója}
%alakpár. Gyakorlati megfontolásból a \a{gó} \emph{-ja} képzős
%alakját hibásnak jelzi a Magyar Ispell: \w{gója}, \w{góját}.
%A Helyes-e?-ben a \w{gó} tő sincs elfogadva.)}
%A kék színnel jelölt szövegdoboz valamilyen plusz szolgáltatást,
%általában javítási képességet ír le, a szürke a
%helyesírás-ellenőrzéshez nem kapcsolódó pontokat jelöli.

\pagebreak
\begin{longtable}[l]{p{7.8cm}p{7.8cm}}
\hline \hfil Magyar Ispell
\hfil & \hfil MS Office XP helyesírás-ellenőrzője \hfil \\ \hline \endhead \hline \endfoot \\

\title{Általános tudnivalók}
\addcontentsline{toc}{section}{Általános tudnivalók}

\s
{green}
{1.\footnote{\ldots{}A helyesírás
valamely nyelv írásának közmegállapodáson alapuló és közérdekből
szabályozott eljárásmódja, illetőleg az ezt tükröző, rögzítő és irányító
szabályrendszer\ldots}
Nyitott forráskódú program, amely a közérdeknek megfelelően áttekinthető:
folyamatosan dokumentált, ellenőrzött, javított és szabadon hozzáférhető
mindenki számára.}
{yellow}
{Zárt, dokumentáció nélkül megjelenő
helyesírás-ellenőrző, amely legalább az MS Office 97 óta bármiféle javítás
nélkül jelenik meg az MS Office különféle változataiban. Linuxos változata
\emph{Mspell} néven otthoni felhasználásra ingyenes.
}

\empty{2.}

\title{A betűk}

\s
{green}
{3. A magyar ábécé minden betűjét önálló alakban helyesnek fogadja el.}
{lightred}
{A magyar ábécé többjegyű betűit hibásnak jelzi: \emph{\w{dzs} betű}.
Az ékezetes betűk toldalékolt alakjait hibásnak jelzi: \emph{hosszú \w{ú-val}}.}


\empty{4--5.}

\s
{yellow}
{6. A stilisztikai célból megengedett
magánhangzó-halmozást nem fogadja el helyesnek: \w{Één?},
\w{Neeem.} Kivéve \a{góóól}, negyven \a{ó}-ig bezárólag.}
{yellow}
{Hasonlóan, de a \w{góóól} alakot hibásnak jelzi.}

\s
{green}
{7.b. A mássalhangzók hosszúságát betűkettőzéssel jelöli.
A többjegyű mássalhangzókat csak az összetett szavak tagjainak határán kettőzi
teljesen, egyéb esetekben egyszerűsítve kettőzi őket.}
{lightred}
{Hiba a \a{-szerű} utótag
összetett szóként való kezelése: \b{észszerű},
\b{vadászszerű}, a helyes alakok (kivétel \a{ésszerű})
hibásnak jelzése: \w{mésszerű}, \w{viasszerű}. L. még 62.}

\s
{green}
{8. Nagybetűs formák esetében
a megengedettek vannak elfogadva: \a{csibe}, \a{Csibe}, \a{CSIBE},
de \emph{\b{\uwave{CSibe}}}, \emph{\b{\uwave{cSiBe}}}.
Gyakorlati megfontolásból
csak nagybetűvel (tulajdonnévként) kerül elfogadásra a \a{kis}
legtöbb toldalékolt alakja (például \w{kissel},
de \a{kissé} jó. }
{green}
{Hasonlóan jól működik. (L. még 164.)}

\s
{yellow}
{9. A stilisztikai célból megengedett
mássalhangzó-halmozást nem fogadja el helyesnek: \w{Nnem!},
\w{Terringettét!}}
{yellow}
{Hasonlóan.}

\empty{10--12.}

\s
{green}{13. A program kezeli az UNICODE
karaktereket. Alapértelmezett az ISO 8859--2
karakterkészlet:\footnote{A közép-európai nyelvek közül az albán, bosnyák, cseh, 
horvát, lengyel, német, magyar, román, szerb, szlovák, szlovén nyelvek
ábécéjének betűit tartalmazza, továbbá az angol, finn és ír ábécé betűit is
lefedi. A Windows operációs rendszer ún. közép-európai kódtáblája (Windows-1250)
egy kibővített ISO 8859--2 kiosztásnak felel meg.} 
\a{Dvořák}, \a{Händel}, \a{Škoda},
\a{złoty}\ldots{} stb. Az UNICODE karakterek ISO 8859--1 esetén
HTML karakterentitással, egyébként pedig decimális
karakterkóddal adhatók meg: \emph{Moli\&egrave;re} (Moli\`ere),
\emph{Duna\&ndash;Tisza köze} (Duna--Tisza köze).}
{lightred}
{Csak a magyar ábécét, és a
\a{q}, \a{w}, \a{x}, \a{y} betűket kezeli a szótár:
\w{Škoda}, \w{Moli\`ere} stb. 
A \b{Skoda}, vagy a \b{zloty} alakok elfogadása
hiba az UNICODE alapú, vagy ISO-8859--2-t támogató
szövegszerkesztők -- mint az OpenOffice.org, illetve a Microsoft Office --
esetében.\footnote{Az OpenOffice.org
teljes UNICODE támogatással rendelkezik, lehetőséget nyújt még az ázsiai
írások kezelésére is. A Helyes-e? és az OpenOffice.org egybeépítésére
a kereskedelmi forgalomban kapható Magyar Office mutat példát.}
}

\s
{cyan}
{Az idegen ékezetek esetében is javítva vannak az ékezet nélküli
változatok: \a{Skoda}\,$\rightarrow$\emph{Škoda}, \a{Dvorák}, sőt
\a{Dvorak}\,$\rightarrow$\emph{Dvořák}.
%\footnote{A
%hatékonysági korlátok miatt mindez nem az ACCENT kombinatorikus ékezetesítéssel,
%hanem a REP cseretáblázat kézzel való bővítésével történt. A teljes igényű
%ellenőrzéshez és javításhoz vagy kiegészítő, más karakter kódolású szótárakra, és ezek
%egyidejű kezelésére van szükség, vagy a REP cseretáblázat hatékonyabb
%kezelésére, amely magában foglalja az idegen ékezetes szavak ékezet nélküli
%megfelelőinek automatikus előállítását és cseretáblázatba való elhelyezését.
%A több különböző szótár egyidejű kezelése már részben megvalósult az OpenOffice.org
%szövegszerkesztőben, ahol lehetőség van több magyar, akár eltérő kódolású szótár
%egyidejű kezelésére is. Sajnos azonban a \emph{\a{Moli\`ere-ről}}, és hasonló
%alakok számára ez a módszer sem jelent megoldást, tehát valóban szükség van
%az egy szótáron belüli UNICODE kezelésre.}
}
{white}
{}

\empty{14--16.}

\title{A kiejtés szerinti írásmód}


\s
{yellow}
{17--19. Ragaszkodik a helyesírás által
rögzített köznyelvi formához. A leglényegesebb hibák (tipikus tévesztések;
a hasonló, de nem rögzített helyesírású alakok előfordulása; hibás, de
a szótárba mégis felvett szóösszetételek) alapján ellenőrizve lett a
szókincs. A teljes szókincs lektorálására azonban még nem került sor.}
{lightred}
{Számos, a helyesírásban nem rögzített
beszélt nyelvi ingadozást helyesnek fogad el. Ilyen, később
nem említett példák: \b{lyány} (\a{lány}),
\b{siserehad} (\a{siserahad}), \b{bédekker}
(\a{bedekker}), \b{ekhó} (\a{echó}), \b{trikolor} (\a{trikolór}),
stb. Valahol csak a rossz alakot engedi: \b{puzon} (\w{puzón}). L. még 28.
\\ Egyéb probléma, hogy nem létező szavakat, illetve
nem szótári töveket tartalmaz a szótár: \b{ál}, \b{fesz};
\b{sepr}, \b{sodr}; \b{süv} \b{veé} stb.
\\ Nem fogad el számos mai magyar szót, különösen
az új, szaknyelvi és földrajzi elnevezéseket: \w{karalábé},
\w{klónozás}, \w{ombudsman}, \w{szafari},
\w{helló} vagy \w{heló} (helyette \b{hello}) stb.
\\ A szótár pontatlan, illetve elnagyolt besorolásokat
tartalmaz. Például hangrendi tévedések: \b{tettról}
(\a{tettről}), \b{vésznak} (\a{vésznek}) stb. Elnagyolások:
\b{fiája}, \b{híjája}, \b{aljájáról} (többszörös birtokos személyragok) stb.; \b{történed},
\b{halja}, \b{évődi} stb. (a nem tárgyas igék kezelésének hiánya).}

\s
{cyan}
{Szaknyelvi modulokat (matematika, informatika stb.)
tartalmaz több száz szakszóval. A történelmi és a mai magyar
helységneveket is tartalmazza, ami a közigazgatási, valamint
földrajz és történelem tárgyú munkák ellenőrzése céljából
nélkülözhetetlen.
A szókincs a felhasználók által is bővíthető.}
{lightred}
{
Nincsenek szaknyelvi modulok: \w{elempár}, \w{alterek},
\w{variancia}, \w{bijektív}, \w{attraktor};
\w{flopi}, \w{mikrocsip}, \w{élkiemelés}, \w{interpreter},
\w{alhálózat}. A városi rangú települések neve közül is sokat nem ismer:
\w{Gyomaendrőd}, \w{Szécsény}, \w{Vásárosnamény}, \w{Zalaszentgrót} stb.
\\ Nincs lehetőség a szókincs bővítésére, csak egyes szóalakok felvételére.}

\s
{green}
{20--21. Az \a{i}, \a{u}, \a{ü} és ezek hosszú változatát mindig ugyanolyan
formában tartalmazó szavak írásában pontos.\footnote{Az alapszókincsre
vonatkozik a pontosság. A magyar helységneveket tartalmazó modul
sok kettős alakot tartalmaz még.}}
{lightred}
{Számos hibás alakot is elfogad:
\b{egyivású} (\a{egyívású}), 
\b{elitélt} (\a{elítélt}), 
\b{hivő} (\a{hívő}),
\b{izület} (\a{ízület});
\b{neutrinó} (\a{neutrínó}),
\b{szivató} (\a{szívató}),
\b{telitett} (\a{telített}),
\b{tikfa} (\a{tíkfa});
\b{kompatíbilis} (\a{kompatibilis}),
\b{színtű} (\a{szintű});
\b{labdarugócsapat} (\a{labdarúgócsapat});
\b{mélyfurás} (\a{mélyfúrás});
\b{ramazuri} (\a{ramazúri});
\b{harangbugás} (\a{harangbúgás});
\b{lámpabúra} (\a{lámpabura});
\b{kűzdőtér} (\a{küzdőtér}) stb.}

\s
{cyan}
{
Gyakorisági megfontolásból sok
alakpár ritkán használt tagja le van tiltva: \w{csábit}, \w{öblit};
\w{visszairt} stb.\footnote{L. a \url{szotar/kivetelek/ragozatlan/i}
állományt a Magyar Ispell forrásában.}}
{green}
{
Egy-két alak van letiltva: \w{tanit}, de például a \a{csábit},
\a{öblit}, \a{visszairt} nem.
}

\s
{green}
{22. Az \a{i}, illetve \a{í} végű szavak esetében
nincs ismert hiba.
}
{lightred}
{Sok ilyen főnév
toldalékolása hibás, illetve hiányos: \w{kocsijai},
\b{kocsiija}, \b{kocsiijai}.
Hasonlóan \w{zsenijei}, \w{mozijai},
\w{masnijai} stb. A \a{hí} toldalékolása
is hibás: \w{híttak}, \b{hittak} stb. 
}

\s
{green}
{23--24. Az \a{u}, \a{ú}, \a{ü}, \a{ű} végű szavak esetében
nincs ismert hiba.}
{lightred}
{Elfogadásra kerülnek a
\b{hiu}, \b{savanyu}, továbbá \b{átlagáru}, \b{vasárú} stb. alakok.
A \a{hú}, \a{hű} indulatszavak mellett helyesli a \w{hu},
és \w{hü} alakokat.}

\s
{green}
{25.
Kezeli a \a{csepp} és \a{csend} alakváltozatai közötti különbséget.
}
{yellow}
{A nem használt alakváltozatokat
is elfogadja: például \b{csöppkő}, \b{csöndőr}.
}

\empty{26.}

\s
{green}
{27. Az \a{í}, \a{ú}, \a{ű} szabálytalan és szabályos
váltakozását mutató szavak írásában pontos.
Gyakorlati megfontolásból viszont nem fogadja el
a ritka \w{tűntet} és \w{szűntet} alakokat (a \a{tűnt},
és \a{szűnt} tárgyragos változatát).
}
{lightred}
{a) A szabálytalan váltakozást mutató szavak
esetében téves alakokat is elfogad:
\b{hivat} (\a{hívat}), \b{hivató} (\a{hívató});
\b{irogat} (\a{írogat}), \b{irogató} (\a{írogató}), \b{irogatás} (\a{írogatás});
\b{bíztat} (\a{biztat}), \b{bíztató} (\a{biztató}), \b{bíztatás} (\a{biztat});
\b{szűntető} (\a{szüntető}), \b{szűntetés} (\a{szüntetés});
\b{tűntető} (\a{tüntető}), \b{tűntetés} (\a{tüntetés}). \\
b) A szabályos váltakozást mutató szavak esetében téves
a \a{fűz} toldalékolása: \b{fűzek} (\a{füzek}),
\b{fűzes} (\a{füzes}) stb. Súlyos hiba mutatkozik még az ilyen
szavak I/3. birtokos személyragos plusz birtokjeles toldalékolásában:
\b{vízéé} (\a{vizéé}, \a{vízé}), \b{tűzéé} (\a{tüzéé}, \a{tűzé}),
\b{kézéé} (\a{kezéé}, \a{kézé}) stb.
Hasonlóan a \a{-nyi} képzős alakok esetében:
\b{viznyi} (\a{víznyi}), \b{tüznyi} (\a{tűznyi}), \b{keznyi} (\a{kéznyi}) stb.}

\s
{green}
{28. A közkeletű idegen szavak magánhangzó-váltakozásában nincs ismert hiba.}
{lightred}
{A \a{kultúra} szó származtatott alakjait
helytelenül kezeli: \b{kultúrál} (\a{kulturál}),
\b{kultúrált} (\a{kulturált}), \b{kultúrálatlan} (\a{kulturálatlan}),
\b{kultúrálódik} (\a{kulturálódik}).
Számos egyéb esetben is elfogadja még a nem rögzített írásmódú alakokat:
\b{analfabétizmus} (\a{analfabetizmus}), \b{polémikus} (\a{polemikus}),
\b{aszkétikus} (\a{aszketikus}), \b{szférikus} (\a{szferikus}) stb.
}

\s
{green}
{29. A szó végi \a{a}, és \a{e} a megfelelő toldalékolás során válik hosszúvá.}
{lightred}
{A \emph{-szerű} toldalék előtt is megnyúlik a magánhangzó minden esetben:
\b{kecskészerű} (\a{kecskeszerű}), 
\b{allegrószerű} (\a{allegroszerű}),
\b{mangrovészerű} (\a{mangroveszerű}).}

\s
{green}
{30. Nem követ el hibát a kettős alakok megengedésében.}
{lightred}
{Nem létező kettős alakokat is megenged a \a{-nyi} toldalék
esetében: \b{borjnyi} (\a{borjúnyi});
\b{tetvnyi} (\a{tetűnyi}); \b{bokrnyi}
(\a{bokornyi}) stb. \\ Az előzőhöz hasonló probléma, hogy birtokjel is
járulhat a nem szótári tőhöz:
\b{bokré} (\a{bokoré}, \a{bokráé}),
\b{terhé} (\a{teheré}, \a{terhéé}) stb.}

\emptyx{green}{31.}

\s
{green}
{32. A \a{metsz}, \a{látszik}, \a{tetszik} toldalékolt alakjai is helyesek.}
{lightred}
{A \w{mesd}; \w{lássék}, \w{tessem}, \w{tessek} stb. alakokat
hibásnak jelzi.}

\emptyx{green}{33--40.}

\s
{green}
{41. A műveltető képző kezelése helyes. \\
A régies \a{-atik}, \a{-etik}, \a{-tatik}, \a{-tetik} szenvedő igeképző
kezelése csak a mai köznyelvben
még megtalálható alakokra szorítkozik: \a{megadatik},
de \w{oltatik}.}
{lightred}
{Az \a{-at}, \a{-et} és a \a{-tat}, \a{-tet} szabadon járul a tőhöz:
\b{bonttat} (\a{bontat}); \b{kértet} (\a{kéret});
\b{nyitat} (\a{nyittat}), \b{oktatatás}\footnote{Érdekes módon
ez egy gyakori tévedés is egyben, aminek hátterében az írás során
megismételt szótag áll. Az ilyen hibákat felismeri és javítja
a Magyar Ispell.} (\a{oktattatás}), \b{történtetnek}, \b{történessél}
(\a{történik}$+$\a{-et}) stb. \\
Általánosan kezeli a szenvedő igeképzőt,
de a műveltető képzőhöz hasonlóan hibásan: \a{oldtatik},
de \b{olttatik}, \w{oltatik}).}

\s
{green}
{42. A \a{-val}, \a{-vel} kezelése jó.}
{yellow}
{Szintén, egy hibával: \b{mohval} (\a{mohhal}).}

\emptyx
{green}
{43--44.}

\s
{green}
{45. A toldalékolásbeli kettősséget jól kezeli,
kivéve gyakorlati okból az elavuló
formákat: \w{szeretniök}, \w{látniok}.}
{green}
{Hasonlóan.}

\s
{green}
{46. Az elhomályosult képzésű szóelemek írásában nincs ismert hiba.}
{lightred}
{Nem ismeri a \a{lélegez} ige egyes alakváltozatait: \w{belégzem},
\w{kilégzel} stb. \\
Több \emph{-szt} végű igét sem kezel:
\w{aggaszt},
\w{apaszt},
\w{áraszt},
\w{ébreszt},
\w{fagyaszt},
\w{függeszt},
\w{horgaszt},
\w{kepeszt},
\w{lehiggaszt},
\w{sorvaszt},
\w{szállaszt}.}

\emptyx{green}{47--48.}

\title{A szóelemző írásmód}

\emptyx{green}{49--61.}

\s
{cyan}
{Több tucat --, a beszélt nyelvi ingadozások,
és hasonulások miatt elkövetett, több karaktert is érintő -- tipikus tévesztést
javít:
\a{elösször}\,$\rightarrow$\a{először},
\a{szöllő}\,$\rightarrow$\a{szőlő},
\a{hüttő}\,$\rightarrow$\a{hűtő},
\a{dijjas}\,$\rightarrow$\a{díjas},
\a{csug}\,$\rightarrow$\a{csukd},
\a{baráccság}\,$\rightarrow$\a{barátság},
\a{licensz}\,$\rightarrow$\a{licenc} stb.\footnote{L. cseretáblázat leírását.}
}
{yellow}
{A \emph{ly}/\emph{j},
\emph{ts}/\emph{cs},
\emph{ggy}/\emph{dj},
\emph{ggy}/\emph{gyj},
\emph{nny}/\emph{nj} több karaktert érintő tévesztéseket javítja csak.
}

\s
{green}
{62. Az azonos hosszú és rövid mássalhangzók találkozását jól kezeli.}
{lightred}
{A \a{-szerű} és a \a{-féle} képzőszerű
utótag esetében hibás alakokat is elfogad: \b{dzsesszszerű}
(\a{dzsessz-szerű}), \b{puffféle} (\a{puff-féle}). \\
A toldalékolásnál helyenként a nem egyszerűsített, hibás alakokat
is elfogadja:
\b{frissség}, \b{fessség}; \b{feddd};
\b{izzzak}, \b{izzzam}\ldots{}\footnote{\ldots{}\b{izzzalak},
\b{izzz},
\b{izzzál},
\b{izzzad},
\b{izzzon},
\b{izzza},
\b{izzzunk},
\b{izzzuk},
\b{izzzatok},
\b{izzzátok},
\b{izzzanak},
\b{izzzák}.} Egyéb példa még a \b{veddd}. A \a{tovább} igekötő esetében
csak a hibás alakot fogadja el: \b{továbbbotorkál} (\w{tovább-botorkál}),
\b{továbbburjánzik} (\w{tovább-burjánzik}) stb.
}

\s
{green}
{63. A \a{d} végű igék ragozása jó.}
{lightred}
{A \a{d} végű igék módjára, de hibásan ragozza az
\a{r}, és \a{j} végű igéket is: \b{írdd} (\a{írd}, \a{írjad}), \b{várdd} (\a{várd},
\a{várjad}); \b{fújdd} (\a{fújd}, \a{fújjad}), \b{vájdd} (\a{vájd}, \a{vájjad}).
}

\s
{green}
{64. Helyesen kezeli az \a{-l} végű igéket.}
{yellow}
{Nem ismeri a \emph{múlik} ige kötőhangzót tartalmazó
alakjait: \w{múlanak}, \w{múlanának}. Elfogadja viszont
az \a{-at}, és \a{-tat} képzős alakjait: \b{múlat}, \b{múltass} (\a{mulat},
\a{mulaszt}). \\
Nem ismeri a \emph{telik} ige hosszú \a{l}-lel írt
alakjait: \w{tellene}, \w{tellenek} stb.}

\s
{green}
{65. Helyesen kezeli az egyalakú \a{-ll} végű igéket.}
{yellow}
{Elfogadja a következő alakot: \b{álnak} (\a{állnak}) stb.\footnote{
Ennek hátterében a hibás \b{ál} (\a{ál-}) tő felvétele \a{áll}.}
}

\emptyx{green}{66.}

\s
{green}
{67. A hosszú magánhangzóra végződő igék ragozásában közel
megegyezik a Helyes-e?-vel. A \a{nőszik} igéből képzett
\w{nőssz} alak a jóval gyakoribb \a{nősz} ellenőrzése
miatt nincs elfogadva.
}
{yellow}
{Hiányzó tő: \w{nyí}, \w{nyíjatok}.
Hibás alak: \b{hittak} (\w{híttak}). Nem fogadja el a régies \a{hí},
és \a{fú} egyes toldalékolt alakjait: (\w{fúvok}, \w{hínak}
stb).
}

\s
{green}
{69. A \a{jön} igét helyesen kezeli.}
{yellow}
{Hiányzó alak: \w{jőnek}.}

\s
{green}
{70. Az \a{egy} számnevet helyesen kezeli.}
{yellow}
{Hibás alak: \b{egyűvé} (\a{együvé}).}

\s
{green}
{72. A \a{kis} fokozott alakjaiban a hosszan ejtett
\a{s}-t írásban nem tükrözi.}
{yellow}
{Hasonlóan, de elfogad egy-két zavaró
alakot hibás szóösszetétel-képzés révén:
\b{kisseb} (\a{kisebbség}), \b{kissebség} (\a{kisebbség}).}

\s
{green}
{72. Az \a{-s} végű melléknevekben a hosszan ejtett
\a{s}-t helyesen kezeli.}
{yellow}
{Hasonlóan, de elfogad egy zavaró
alakot hibás szóösszetétel-képzés révén:
\b{erősseb} (\a{erősebb}).}

\emptyx{green}{73.}

\s
{cyan}
{
Helyes javaslatot tesz a \a{belölle},
\a{elölle} hibás alakokra: \a{belőle},
\a{előle}.}
{yellow}
{Nem tesz javaslatot}

\s
{green}
{74. Az \a{új} és \a{ujj} meg van különböztetve.}
{yellow}
{Az \a{ujj} \a{-ság} képzőt is kaphat:
\b{ujjság} (\a{újság}). Az \a{új} a
főnevekhez hasonlóan bármely szóösszetételben
szerepelhet: \b{újperc} (\a{ujjperc}).
}

% pechkel?

\s
{green}
{75. A \a{h} végű szavak írása pontos.}
{yellow}
{Szintén, kivéve \b{mohval} (\a{mohhal}).
}

\empty{76.}

\emptyx{green}{77.}

\s
{green}
{78. A rövid magánhangzó $+$ \a{t} végű igék, illetve a
\a{lát} és \a{bocsát} igék kezelése helyes.}
{lightred}
{A hibás \b{lássd} (\a{lásd}), \b{bocsássd} (\a{bocsásd}), \b{hallgassd} (\a{hallgasd}),
\b{kössd} (\a{kösd}), \b{nyissd} (\a{nyisd}) stb. alakokat elfogadja.}

\s
{green}
{79. Az \a{st} és \a{szt} végű igék kezelése helyes.}
{lightred}
{Az \a{st} végű igénél (\a{fest}) elfogadja a \b{fessd} (\a{fesd})
alakot.}

\s
{green}
{80--81. Az \a{s}, \a{sz}, \a{z} és \a{dz} végű igék ragozása helyes.}
{lightred}
{Az \a{s}, \a{sz} és \a{dz} végű igék kezelése hibás: 1. \b{olvassd} (\a{olvasd}),
\b{késsd} (\a{késd}) stb.; 2. \b{mászza} (\a{mássza}), \b{játszzák} (\a{játsszák})
stb.; 3. \w{tessem}, \w{tessél}, \w{lássék}, \w{mesd}. 4. \w{edzd}, \w{pedzd}.
}

\s
{green}
{82. A kivételes ragozású \a{sz} végű igék kezelése helyes.}
{yellow}
{A \a{tesz}, \a{visz}, \a{eszik} kezelése hibás: \b{tesszél} (\a{tegyél}),
\b{visszétek} (\a{vigyétek}), \b{esszünk} (\a{együnk}) stb.
}

\s
{green}
{83. A \a{-val}, \a{-vel} toldalékolás jó.}
{green}
{Hasonlóan, a \b{mohval}, és a betűk toldalékolása (l. 3.)
kivételével.
}

\emptyx{green}{84--85.}

\title{A hagyományos írásmód}

\empty{86.}

\emptyx{green}{87--88.}

\s
{green}
{89. Az \a{ly} írásában pontos.}
{lightred}
{Csak a hibás alakot fogadja a \w{csevej}
esetében: \b{csevely}. \\
Számos \a{ly}-t tartalmazó szó
esetében a hibás alakot is elfogadja:
\b{estéji}, \b{karvaj}, \b{szeméji},
\b{ünnepéjes}; \b{bolyár}, \b{bolytár},
\b{súlytó} stb. 
}

\s
{cyan}
{
A keretrendszer lehetőséget ad a zárt \a{\"e} hangokat jelölő
helyesírási szótár létrehozására a magyar nyelv zárt \a{\"e} szótára
alapján.
}
{white}
{
}

\empty{green}{91.}

\title{Az egyszerűsítő írásmód}

\empty{92--93.}


\s
{green}
{94. Jól kezeli az egyszerűsítő írásmódot. Az ezekben elkövetett
hibákra javaslatot tesz (például Gaussal$\rightarrow$Gauss-szal.)}
{lightred}
{Akadnak típushibák, például \w{vadásszerű}, \b{vadászszerű}, vagy 
\w{puff-féle} és \b{puffféle}. A helyes kötőjeles alakot
nem javasolja.}

\title{A különírás és az egybeírás}

\s
{yellow}
{95. Fokozott hangsúly az egybe- és különírási szabályok betartásán.
A helyesírási szótár szintjén kezel nagyon sok szóalakot.
A kötőjeles szavak írásánál viszont még nagyon megengedő.
}
{lightred}
{Alapvető hiányosságok (például nagyon sok esetben
melléknevek és főnevek egybeírásának engedélyezése).
A kötőjeles szavak írásánál is nagyon megengedő.
}

\s
{green}
{96--99. Különírandó szóismétlések írásánál feltünteti a helyes kötőjeles
alakot is.
}
{yellow}
{Nem javasolja a kötőjeles alakot.
}

\s
{green}
{106. Az alanyos kapcsolatok esetében a jelentésváltozással bírókat
engedi egybeírni.
}
{lightred}
{Minden alanyos kapcsolatot elfogad (\b{macskaásított}, \b{egérfutott}.
}

\s
{green}
{107--108. A minőségjelzős kapcsolatok esetében a jelentésváltozással bírókat
engedi csak egybeírni.
}
{lightred}
{Sok minőségjelző esetében minden kapcsolatot elfogad (\b{nagyépület},
\b{butakutya}, \b{szépvers}.
}

\s
{green}
{111. Az önálló alakban nem szereplő minőségjelzős alakokat nem fogadja el,
csak külön felvett összetett szóban, vagy kötőjellel (\a{al-}, \a{ál-}, \a{bel-},
\a{fesz-}).
}
{lightred}
{Több ilyet elfogad: \b{ál}, \b{bel}, \b{fesz}, \b{pót}.
}

\s
{yellow}
{112--114. Az -ó, -ő képzős minőségjelzős összetett alakokat, ha a szó valamely
tagja összetett, szótári tiltások alapján a leggyakoribb esetekben nem engedi egybeírni.
Az egybeírt alakokra helyes javaslatot tesz. A lista nem lehet teljes,
de a leggyakoribb hibás egybeírásokat tartalmazni fogja a Szószablya fejlesztésnek
köszönhetően.
}
{lightred}
{Nagyon sok különírandó alakot helyesen kezel, de így is sok gyakorit enged
egybeírni: \b{ellenőrzőbizottság}, \b{előadókörút}.
}

\s
{lightred}
{115. Az anyagnevek felismerése csak a gyakori esetekre van megoldva.
}
{lightred}
{Hasonlóan.
}

\s
{green}
{116. A hónapok, illetve napok nevét nem engedi egybeírni.
}
{lightred}
{A hónapok, illetve napok nevét engedi egybeírni.
\b{januárhónapban}, \b{hétfőreggel}.
}

%\colorbox{green}{\parbox[t]{7.8cm}{164.
%}} &
%\colorbox{lightred}{\parbox[t]{7.8cm}{
%A nagybetűs írásmódot érintő szabály, egy-két
%kivételtől eltekintve: \b{Rt}. \\
%Az \emph{-s} képzős tulajdonnevek esetében a nagybetűs alakok
%vannak hibásan kikényszerítve:
%\w{trabantos}$\rightarrow$\a{Trabantos}}} \\

\end{longtable}

%\textcolor{blue}{A hibás alakoknál
%helyes javaslatot tesz: \emph{\uwave{baletttáncos}$\rightarrow$balett-táncos},
%\emph{\uwave{dzsesszszak}$\rightarrow$balett-táncos}}

%Az a szó, hogy \uwave{csevely}, hibás.
%Színesen: \textcolor{lightred}{\uwave{\textcolor{black}{csevely}}}.


